\section{Conclusion}

\begin{center}
\begin{tabular}{p{0.37\textwidth}p{0.24\textwidth}p{0.27\textwidth}}
\hline
Objective & Polynomial time & Quasipolynomial time\\
\hline
Maximum depth, binary answers &
$2+\gamma=4.46496526$ & $2+\varepsilon$\\[3pt]
Maximum depth, on-line answers &
$4+2\gamma=8.92993052$ & $4+\varepsilon$\\[3pt]
Expected depth, binary answers and arbitrary priors &
$\rho_{\rm two}<5.182976$ & $2+\varepsilon$\\[3pt]
Expected depth, on-line answers and arbitrary priors &
$2\rho_{\rm two}<10.365952$ & $4+\varepsilon$\\
\hline
\end{tabular}
\end{center}

The common-witness normalization is the structural step shared by
both algorithms. In a flat state, it turns each supplied test into
an opposite cap and represents that cap by at most two adjacent
vertices of a convex hull. For each hypothesis, the positive hull
answers form a circular interval. This representation does not
enlarge the query family: it supplies an interval view of the tests
that were already given.

The two objectives then use that structure differently. For maximum
depth, the algorithm schedules discovery around the actual heights
of its recursively built children. The comparison path in an
optimal tree supplies the deadline credits, while adjacent support
blocks control discovery delay; route-specific endpoint searches
preserve those credits. Exact Cartesian-tree scheduling with rank
subsets handles logarithmically many discovery slots in polynomial
time. A global slot cap and the priority fallback support simultaneous
small-optimum and universal bounds, yielding factor $4.464966$.
For expected depth, the algorithm
learns each full hull pattern before recursing. Queries that become
unnecessary after restriction pay for discovery, and the entropy
still remaining after discovery pays for endpoint information.
Exact integer-budget coverage applied to adjacent hull blocks,
joint discovery and endpoint accounting, and the zero-pattern waiting
time sharpen the polynomial factor. Retaining an orientation credit
through the first two batches' coverage guarantees lowers it to $5.283310$.
Rational choices between stopping and querying both endpoints retain
the second answer's conditional information and lower it further to
$5.272565$. Coarsening the transcript preserves the individual cap
bound without charging the information twice. Changing the initial
coverage budget and merging selected first phases lowers the factor to
$5.197569$. Modifying both initial budgets lowers it further to
$5.182976$. Signed first-cover corrections and a separate all-negative
waiting calculation account for every effect of these changed schedules.
In both analyses, the local credits
are decreases of quantities that telescope over the recursive tree.

The quasipolynomial algorithms retain a logarithmic number of schedule
grades, giving direct factors below $4.464966$ and $3.7187$.
Their additive bounds have leading coefficient arbitrarily close to
two. A logarithmic-depth lookahead dynamic program then removes the
count and entropy surcharge from the ratio, proving factors
$2+\varepsilon$ for every fixed $\varepsilon>0$.
The on-line-answer entries follow from the same factor-two simulation,
using the binary guarantee with parameter $\varepsilon/2$.
These results have quasipolynomial running time.

The small-cap results place additional limits on potential hardness
constructions. With a distinguished all-miss hypothesis and fixed
cap size, maximum depth is exactly solvable in polynomial time.
The weighted singleton/pair path-and-cycle subclass with all
singleton tests has an EPTAS for expected depth. Thus the current
weighted pair-path reduction cannot yield APX-hardness unless
$\mathrm P=\mathrm{NP}$.
Uniform fixed-size caps in convex position with an all-miss hypothesis
also admit a PTAS, without requiring singleton tests. Clearing short
gaps permits exact optimization of the residual block interleavings;
the prefix cost is a small fraction of the full optimum. Constant-factor
hardness must therefore escape this additional restricted class.

The hardness reductions show that this geometry does not make exact
optimization easy: even points on a parabola with only supplied
binary line queries encode strongly NP-complete 3-Partition
instances. The expected-depth result holds for both arbitrary positive
priors and uniform priors: a uniform cloud lift adds a fixed internal
cost to the weighted optimum. Constant-factor inapproximability
remains open. Whether the polynomial factors can be lowered or
the quasipolynomial factors can be brought below two remains open.
